\documentclass[11pt]{article}
\usepackage[a4paper,margin=1in]{geometry}
\usepackage{amsmath,amssymb}
\usepackage{hyperref}
\title{Vectorized Miller Planes and Directions in PyTex}
\author{PyTex Contributors}
\date{\today}

\begin{document}
\maketitle

\section{Scope}

This note records the scientific contract for the first-class PyTex Miller surface:

\begin{itemize}
  \item scalar \texttt{MillerPlane} and \texttt{MillerDirection},
  \item batch \texttt{MillerPlaneSet} and \texttt{MillerDirectionSet},
  \item vectorized conversion between integer indices and cartesian direct or reciprocal vectors,
  \item family canonicalization, antipodal handling, and symmetry-equivalent expansion.
\end{itemize}

The design goal is to preserve exact integer input while making family operations explicit and
vectorized.

\section{Stored Indices And Bases}

PyTex stores plane indices as integer reciprocal-basis coordinates $(h, k, l)$ and direction
indices as integer direct-basis coordinates $(u, v, w)$. The basis matrices are:

\[
\mathbf{A} = [\mathbf{a}\ \mathbf{b}\ \mathbf{c}],
\qquad
\mathbf{B} = \mathbf{A}^{-\mathsf{T}}.
\]

The reciprocal convention is the crystallographer convention without a $2\pi$ factor, so

\[
\mathbf{a}_i^* \cdot \mathbf{a}_j = \delta_{ij}.
\]

For a direction index row vector $\mathbf{u}$ and plane index row vector $\mathbf{h}$, the
cartesian vectors are

\[
\mathbf{r} = \mathbf{u}\mathbf{A}^{\mathsf{T}},
\qquad
\mathbf{g} = \mathbf{h}\mathbf{B}^{\mathsf{T}}.
\]

\section{Normals, d-Spacings, And Angles}

Plane normals are the normalized reciprocal vectors:

\[
\hat{\mathbf{n}} = \frac{\mathbf{g}}{\lVert \mathbf{g} \rVert}.
\]

The interplanar spacing is

\[
d = \frac{1}{\lVert \mathbf{g} \rVert},
\]

reported by PyTex in angstroms.

Direction unit vectors are

\[
\hat{\mathbf{r}} = \frac{\mathbf{r}}{\lVert \mathbf{r} \rVert}.
\]

For angle calculations, PyTex uses

\[
\theta = \arccos\!\left(\operatorname{clip}(\hat{\mathbf{x}} \cdot \hat{\mathbf{y}}, -1, 1)\right),
\]

with antipodal handling applied by replacing the dot product with its absolute value where the
family semantics require it.

\section{Family Canonicalization}

Exact indices are preserved at construction time. Family operations are explicit:

\begin{itemize}
  \item row-wise GCD reduction,
  \item deterministic sign canonicalization by requiring the first non-zero component to be positive,
  \item antipodal family keys given by reduced-and-sign-canonicalized indices.
\end{itemize}

This separation matters because $(2, 0, 0)$ and $(1, 0, 0)$ do not carry the same spacing even
though they belong to the same geometric family.

\section{Hexagonal Four-Index Conversions}

For hexagonal axes, PyTex exposes vectorized integer conversions:

\[
(h, k, l) \leftrightarrow (h, k, i, l), \qquad i = -(h + k),
\]

\[
(u, v, w) \leftrightarrow (U, V, T, W), \qquad U + V + T = 0.
\]

The implementation avoids per-row rational arithmetic objects and performs the batch conversion with
integer formulas plus row-wise GCD reduction.

\section{Symmetry Expansion}

Symmetry operators act on non-normalized cartesian lattice vectors before conversion back to integer
indices:

\[
\mathbf{g}' = \mathbf{R}\mathbf{g}, \qquad
\mathbf{h}' = \operatorname{round}\!\left(\mathbf{g}'\mathbf{B}^{- \mathsf{T}}\right),
\]

\[
\mathbf{r}' = \mathbf{R}\mathbf{r}, \qquad
\mathbf{u}' = \operatorname{round}\!\left(\mathbf{r}'\mathbf{A}^{- \mathsf{T}}\right).
\]

PyTex rejects the conversion if the recovered coordinates are not within tolerance of an integer
index. Family reduction and optional antipodal canonicalization are applied after the integer
recovery step.

\section{Projection Onto Planes}

Direction projection onto a plane uses unit cartesian direction vectors and unit plane normals:

\[
\mathbf{p} = \hat{\mathbf{r}} - (\hat{\mathbf{r}} \cdot \hat{\mathbf{n}})\hat{\mathbf{n}}.
\]

Degenerate rows with $\lVert \mathbf{p} \rVert \approx 0$ are not raised as batch-fatal errors.
PyTex instead returns a boolean degenerate mask alongside the projected vectors.

\section{Normative References}

\begin{thebibliography}{9}
\bibitem{bunge}
H.-J. Bunge,
\textit{Texture Analysis in Materials Science: Mathematical Methods},
Butterworths, 1969.
DOI: \url{https://doi.org/10.1016/C2013-0-11769-2}.

\bibitem{mtex}
R. Hielscher and collaborators,
\textit{MTEX documentation},
\url{https://mtex-toolbox.github.io/}.

\bibitem{orix}
orix documentation,
\url{https://orix.readthedocs.io/}.
\end{thebibliography}

\end{document}
